%% -*- coding: utf-8 -*-
\documentclass{beamer}

%% -*- coding: utf-8 -*-
\usetheme{Boadilla} % default
\useoutertheme{infolines}
\setbeamertemplate{navigation symbols}{} 

\usepackage{etex}
\usepackage{alltt}
\usepackage{pifont}
\usepackage{color}
\usepackage[utf8]{inputenc}
%\usepackage{german}
\usepackage{listings}
\definecolor{keywordcolor}{rgb}{0.7,0.1,0.1}
\definecolor{tacticcolor}{rgb}{0.0,0.1,0.6}
\definecolor{commentcolor}{rgb}{0.4,0.4,0.4}
\definecolor{symbolcolor}{rgb}{0.0,0.1,0.6}
\definecolor{sortcolor}{rgb}{0.1,0.5,0.1}
\definecolor{attributecolor}{rgb}{0.7,0.1,0.1}

\def\lstlanguagefiles{lstlean.tex}
\lstset{language=Lean}
\lstset{sensitive=true}
\usepackage{hyperref}
\hypersetup{colorlinks=true}
\usepackage[final]{pdfpages}
\usepackage{url}
\usepackage{arydshln} % dashed lines
\usepackage{tikz}
\usepackage{mathpartir}

\DeclareUnicodeCharacter{3BB}{\ensuremath{\lambda}}


\newcommand\cmark{\ding{51}}
\newcommand\xmark{\ding{55}}

\newcommand{\nat}{\mathbf{N}}

\usepackage[all]{xy}

%% new arrow tip for xy
\newdir{|>}{!/4.5pt/@{|}*:(1,-.2)@^{>}*:(1,+.2)@_{>}}

\newcommand\cid[1]{\textup{\textbf{#1}}} % class names
\newcommand\kw[1]{\textup{\textbf{#1}}}  % key words
\newcommand\tid[1]{\textup{\textsf{#1}}} % type names
\newcommand\vid[1]{\textup{\texttt{#1}}} % value names
\newcommand\Mid[1]{\textup{\texttt{#1}}} % method names

\newcommand\TODO[1][]{{\color{red}{\textbf{TODO: #1}}}}

\newcommand\String[1]{\texttt{\dq{}#1\dq{}}}

\newcommand\ClassHead[1]{%
  \ensuremath{\begin{array}{|l|}
      \hline
      \cid{#1}
      \\\hline
    \end{array}}}
\newcommand\AbstractClass[2]{%
  \ensuremath{\begin{array}{|l|}
      \hline
      \cid{\textit{#1}}
      \\\hline
      #2
      \hline
    \end{array}}}
\newcommand\Class[2]{%
  \ensuremath{\begin{array}{|l|}
      \hline
      \cid{#1}
      \\\hline
      #2
      \hline
    \end{array}}}
\newcommand\Attribute[3][black]{\textcolor{#1}{\Param{#2}{#3}}\\}
\newcommand\Methods{\hline}
\newcommand\MethodSig[3]{\Mid{#2} (#3): \,\tid{#1}\\}
\newcommand\CtorSig[2]{\Mid{#1} (#2)\\}
\newcommand\AbstractMethodSig[3]{\Mid{\textit{#2}} (#3): \,\tid{#1}\\}
\newcommand\Param[2]{\vid{#2}:~\tid{#1}}

\lstset{%
  frame=single,
  xleftmargin=2pt,
  stepnumber=1,
  numbers=left,
  numbersep=5pt,
  numberstyle=\ttfamily\tiny\color[gray]{0.3},
  belowcaptionskip=\bigskipamount,
  captionpos=b,
  escapeinside={*'}{'*},
  % language=java,
  tabsize=2,
  emphstyle={\bf},
  commentstyle=\mdseries\it,
  stringstyle=\mdseries\rmfamily,
  showspaces=false,
  showtabs=false,
  keywordstyle=\bfseries,
  columns=fullflexible,
  basicstyle=\footnotesize\CodeFont,
  showstringspaces=false,
  morecomment=[l]\%,
  rangeprefix=////,
  includerangemarker=false,
}

\newcommand\CodeFont{\sffamily}

\definecolor{lightred}{rgb}{0.8,0,0}
\definecolor{darkgreen}{rgb}{0,0.5,0}
\definecolor{darkblue}{rgb}{0,0,0.5}

\newcommand\highlight[1]{\textcolor{blue}{\emph{#1}}}
\newcommand\GenClass[2]{\cid{#1}\texttt{<}\cid{#2}\texttt{>}}

\newcommand\Colored[3]{\alt<#1>{\textcolor{#2}{#3}}{#3}}

\newcommand\nt[1]{\ensuremath{\langle#1\rangle}}

\newcommand{\free}{\operatorname{free}}
\newcommand{\bound}{\operatorname{bound}}
\newcommand{\var}{\operatorname{var}}
\newcommand\VSPBLS{\vspace{-\baselineskip}}

\newcommand\IF{\textit{IF}}
\newcommand\TRUE{\textit{TRUE}}
\newcommand\FALSE{\textit{FALSE}}

\newcommand\IFZ{\textit{IF0}}
\newcommand\ZERO{\textit{ZERO}}
\newcommand\SUCC{\textit{SUCC}}
\newcommand\ADD{\textit{ADD}}
\newcommand\SUB{\textit{SUB}}
\newcommand\MULT{\textit{MULT}}
\newcommand\DIV{\textit{DIV}}

\newcommand\PAIR{\textit{PAIR}}
\newcommand\FST{\textit{FST}}
\newcommand\SND{\textit{SND}}

\newcommand\CASE{\textit{CASE}}
\newcommand\LEFT{\textit{LEFT}}
\newcommand\RIGHT{\textit{RIGHT}}

\newcommand\Encode[1]{\lceil#1\rceil}
\newcommand\Reduce{\stackrel\ast\rightarrow_\beta}

\newcommand\Nat{\textit{Nat}}
\newcommand\Bool{\textit{Bool}}
\newcommand\Pair{\textit{Pair}}
\newcommand\Tfun[1]{#1\to}

\newcommand\Tenv{A}
\newcommand\Lam[1]{\lambda#1.}
\newcommand\App[1]{#1\,}
\newcommand\Succ{\textit{SUCC}\,}
\newcommand\Let[2]{\textit{let}\,#1=#2\,\textit{in}\,}

\newcommand\calE{\mathcal{E}}
\newcommand\calU{\mathcal{U}}
\newcommand\calP{\mathcal{P}}
\newcommand\calW{\mathcal{W}}

\newcommand\GEN{\textit{gen}}
\newcommand\EFV[1]{\textit{fv} (#1)}
\newcommand\Dom[1]{\textit{dom} (#1)}

%%% Local Variables: 
%%% mode: latex
%%% TeX-master: nil
%%% End: 

%%% frontmatter
%% -*- coding: utf-8 -*-

\title{Functional Programming / Lean}
\subtitle{Introduction}

\author[Peter Thiemann]{Prof. Dr. Peter Thiemann}
\institute[Univ. Freiburg]{Albert-Ludwigs-Universität Freiburg, Germany}
\date{WS 2026/27}


\subtitle{Interpreters and Monads}

\begin{document}

\begin{frame}
  \titlepage
\end{frame}

%\begin{frame}
%  \frametitle{Outline}
%  \tableofcontents
  % You might wish to add the option [pausesections]
%\end{frame}


\begin{frame}[fragile]
  \frametitle{A simple expression language}
  \framesubtitle{Abstract syntax (AST)}
\begin{block}{Definition}
\begin{lstlisting}
inductive Op where
  | add : Op
  | sub : Op
  | mul : Op
  | div : Op

inductive Term where
  | con : Nat → Term
  | bin : Term → Op → Term → Term
\end{lstlisting}  
\end{block}
\end{frame}             

\begin{frame}[fragile]
  \frametitle{A simple interpreter}
\begin{block}{Evaluation}
\begin{lstlisting}
def opEval : Op → Nat → Nat → Nat
  | .add, x, y => x + y
  | .sub, x, y => x - y
  | .mul, x, y => x * y
  | .div, x, y => x / y

def eval : Term → Nat
  | .con n => n
  | .bin t1 op t2 => opEval op (eval t1) (eval t2)
\end{lstlisting}  
\end{block}
\end{frame}             

\begin{frame}[fragile]
  \frametitle{Extending the interpreter}
  \begin{block}{Possible extensions}
    \begin{itemize}
    \item Error handling
    \item Counting evaluation steps
    \item Variables, state
    \item Output
    \end{itemize}
    \dots{} but without changing the structure of the interpreter!
  \end{block}
\end{frame}

\begin{frame}[fragile]
  \frametitle{Feature: Error handling}
  \begin{itemize}
  \item A computation either returns a value or it fails on the way
  \item Could be expressed with the \lstinline{Maybe} type
    \begin{itemize}
    \item Returning the value \lstinline{v} is \lstinline{Just v}
    \item Signalling an error is \lstinline{Nothing}
    \end{itemize}
  \item But it would be better to have an error message
\begin{lstlisting}
inductive Result α where
  | return : α → Result α
  | raise : String → Result α
\end{lstlisting}
  \item Equivalently we might use the  predefined type \lstinline{Sum α β} (or
    \lstinline{α ⊕ β})
\begin{lstlisting}
abbrev Result' α := String ⊕ α
\end{lstlisting}
  \end{itemize}
\end{frame}

\begin{frame}[fragile]
  \frametitle{Interpreter with error handling}
  \begin{block}{Exception}
\begin{lstlisting}
def evalExc : Term → Exc Nat
  | .con n => Exc.return n
  | .bin t1 op t2 =>
    match evalExc t1, evalExc t2 with
    | Exc.raise msg, _ => Exc.raise msg
    | _, Exc.raise msg => Exc.raise msg
    | Exc.return v1, Exc.return v2 =>
      if op == .div && v2 == 0 then
        Exc.raise "division by zero"
      else
        Exc.return (opEval op v1 v2)
\end{lstlisting}
\end{block}     
\end{frame}

\begin{frame}[fragile]
  \frametitle{Feature: Tracing}
  \begin{itemize}
  \item We would like to produce a trace of the computation
  \item Each step of the computation should be documented
  \item Such a computation returns the computed value \textbf{and} the
    trace
  \item Could be expressed with the following structure type:
\begin{lstlisting}
structure Trace α where
  trace ::
  val : α
  log : String
open Trace

def trace' : α → String → Trace α
  | val, log => { val := val, log := log }
def trace'' : α → String → Trace α
  | val, log => ⟨ val, log ⟩ 
\end{lstlisting}
  \end{itemize}
\end{frame}

\begin{frame}[fragile]
  \frametitle{Interpreter with tracing}
  \begin{block}{Trace}
\begin{lstlisting}
def evalTrace : Term → Trace Nat
  | .con n => trace n s!"con {n}"
  | .bin t1 op t2 =>
    match evalTrace t1, evalTrace t2 with
    | r1, r2 =>
      let v := opEval op r1.val r2.val
      let log := r1.log ++ r2.log ++ s!"evaluated {r1.val} {op} {r2.val} = {v}"
      trace v log
\end{lstlisting}  
\end{block}
\end{frame}

\begin{frame}[fragile]
  \frametitle{Feature: Reduction count}
  \begin{itemize}
  \item We want to keep track of the number of reductions during
    evaluation of an expressions
  \item To this end, a counter value must be passed as a parameter to
    the \lstinline{eval} function, which must return the updated
    counter
  \item A type like the following would do
\begin{lstlisting}
structure Count α where
  count ::
  exCount : Nat → α × Nat
open Count

def tick : Count Unit := ⟨ fun n => ((), n+1) ⟩
\end{lstlisting}
  \end{itemize}
\end{frame}

\begin{frame}[fragile]
  \frametitle{Interpreter with reduction count}\small
  \begin{block}{Count}
\begin{lstlisting}
def evalCount : Term → Count Nat
  | .con n => ⟨ fun c => (n, c) ⟩
  | .bin t1 op t2 =>
    let r1 := evalCount t1
    let r2 := evalCount t2
    ⟨ fun c =>
      let (v1, c1) := r1.exCount c
      let (v2, c2) := r2.exCount c1
      let v := opEval op v1 v2
      (v, c2 + 1) ⟩
\end{lstlisting} 
\end{block}
\end{frame}            

\begin{frame}
  \frametitle{Monads}
  \begin{itemize}
  \item Abstract concept of \emph{computation} behind these
    examples: Monads
  \item Originates in category theory
  \item Mathematicians did not even bother to give it a proper name,
    they called it ``standard construction''
  \item Discovered for computations by Eugenio Moggi
  \item Popularized for programming by Phil Wadler
  \item In Lean: the \lstinline{Monad} type class, organized in two subclasses,
    \lstinline{Bind} and \lstinline{Applicative}.
  \end{itemize}
\end{frame}

\begin{frame}[fragile]
  \frametitle{Building blocks of computations}
  \begin{block}<+->{Composing computations: Bind}
\begin{lstlisting}
class Bind (m : Type → Type) where
  bind : {α β : Type} → m α → (α → m β) → m β
\end{lstlisting}
    \begin{itemize}
    \item Bind \lstinline{(. >>= .)} is for \textbf{composing} computations
    \end{itemize}
  \end{block}
  \begin{block}{Trivial computation: Pure}
\begin{lstlisting}
class Pure (f : Type → Type) where
  pure : {α : Type} → α → f α
\end{lstlisting}
    \begin{itemize}
    \item Pure injects a value in a computation.
    \item The resulting computation does nothing but return that value.
    \end{itemize}
  \end{block}
\end{frame}

\begin{frame}[fragile]
  \frametitle{Monads to the rescue!}
\begin{block}<+->{The type class Monad}
\begin{lstlisting}
class Monad (m : Type → Type) extends Pure m, Bind m
\end{lstlisting}
  Here, \lstinline{m : Type -> Type} is a variable that can stand for
  \texttt{IO}, \texttt{Gen}, and other \textbf{type constructors}. 
\end{block}
\begin{block}<+->{Purpose of the operations}
  \begin{itemize}
  \item Glue operations
    \begin{itemize}
    \item Bind \lstinline{(>>=)} is for \textbf{composing}
      computations\\
      In \lstinline{IO} composing means sequencing
    \item \lstinline{pure} is for injecting values into computations\\
      A trivial computation that does nothing but return the value
    \end{itemize}
  \item Monads live for their other features!
    \begin{itemize}
    \item primitive IO operations
    \item \lstinline{tick}, \lstinline{trace}, etc
    \end{itemize}
  \end{itemize}
\end{block}
\end{frame}     


\begin{frame}[fragile]
  \frametitle{Monadic interpreter}
\begin{exampleblock}<+->{Monadic interpreter}
\begin{lstlisting}
def evalM [Monad m] : Term → m Nat
  | .con n => pure n
  | .bin t1 op t2 =>
    evalM t1 >>= fun v1 =>
    evalM t2 >>= fun v2 =>
    pure (opEval op v1 v2)
\end{lstlisting}  
\end{exampleblock}
  \begin{alertblock}<+->{The identity monad}
\begin{lstlisting}
abbrev Id (α : Type) := α

instance : Monad Id where
  pure x := x
  bind x f := f x
\end{lstlisting}  
\end{alertblock}

\end{frame}             


\begin{frame}[fragile]
  \frametitle{Monadic interpreter with error handling}
  \begin{alertblock}{Exeception}
\begin{lstlisting}
instance Monad Exception where
  return a = Return a
  m >>= f  = case m of 
               Raise  s -> Raise s
               Return v -> f v
  fail s   = Raise s                -- the extra operation
\end{lstlisting}  
\end{alertblock}

\begin{exampleblock}{Interpreter}
\begin{lstlisting}
def evalExcM : Term → Exc Nat
  | .con n => pure n
  | .bin t1 op t2 => do
    let v1 ← evalExcM t1
    let v2 ← evalExcM t2
    if op == .div && v2 == 0 then
      Exc.raise "division by zero"
    else
      pure (opEval op v1 v2)
\end{lstlisting}  
\end{exampleblock}
\end{frame}             



\begin{frame}[fragile]
  \frametitle{A monad for tracing}
  \begin{alertblock}{Trace}
\begin{lstlisting}
instance : Monad Trace where
  pure x := trace x ""
  bind x f :=
    let r := f x.val
    trace r.val (x.log ++ r.log)

def output (msg : String) : Trace Unit :=
  trace () msg
\end{lstlisting}  
\end{alertblock}
\end{frame}

\begin{frame}[fragile]
 \frametitle{Monadic interpreter with tracing}
 \begin{exampleblock}{Evaluation}
\begin{lstlisting}
def evalTraceM : Term → Trace Nat
  | .con n => do
    output s!"con {n}"
    pure n
  | .bin t1 op t2 => do
    let v1 ← evalTraceM t1
    let v2 ← evalTraceM t2
    let v := opEval op v1 v2
    output s!"evaluated {v1} {op} {v2} = {v}"
    pure v
\end{lstlisting}
\end{exampleblock}      
\end{frame}             



\begin{frame}[fragile]
  \frametitle{A monad for counting}
  \begin{alertblock}{Counting}
\begin{lstlisting}
instance : Monad Count where
  pure x := ⟨ fun c => (x, c) ⟩
  bind x f :=
    ⟨ fun c =>
      let (v1, c1) := x.exCount c
      (f v1).exCount c1 ⟩
\end{lstlisting} 
\end{alertblock}
\end{frame}

\begin{frame}[fragile]
  \frametitle{Monadic interpreter with reduction count}
  \framesubtitle{Implementation}
\begin{exampleblock}{Evaluation}
\begin{lstlisting}
def evalCountM : Term → Count Nat
  | .con n => pure n
  | .bin t1 op t2 => do
    let v1 ← evalCountM t1
    let v2 ← evalCountM t2
    tick
    pure (opEval op v1 v2)
\end{lstlisting} 
\end{exampleblock}
\end{frame}            

\begin{frame}[fragile]
  \frametitle{Typical monads}
  \begin{block}{Already used}
  \begin{itemize}       
  \item Identity monad
  \item Exception monad
  \item State monad
  \item Writer monad
  \end{itemize} 
  \end{block}   
\end{frame}

\begin{frame}[fragile]
  \frametitle{Not every type constructor can be a monad}
  \framesubtitle{Monad laws}
  \begin{block}{Assume}
\begin{lstlisting}
x : α            f : α → m β             m : m α
\end{lstlisting}
  \end{block}
  \begin{alertblock}{map lifts identify}
\begin{lstlisting}
∀ {m : Type → Type} {α : Type} (y : f α), map id y = y
\end{lstlisting}
  \end{alertblock}
  \begin{alertblock}{pure is a left unit}
\begin{lstlisting}
∀ {m : Type → Type} {α β : Type} (x : α) (f : α → m β), pure x >>= f = f x
\end{lstlisting}
  \end{alertblock}
  \begin{alertblock}{bind is associative}
\begin{lstlisting}
∀ {m : Type → Type} {α β γ : Type u} (x : m α) (f : α → m β) (g : β → m γ),
  x >>= f >>= g = x >>= fun x => f x >>= g
\end{lstlisting}
  \end{alertblock}
\end{frame}

\begin{frame}[fragile]
  \frametitle{The Option monad}
  \begin{block}{More useful than you think}
                \begin{itemize}         
                        \item Computation that may or may not return a result
                        \item Database queries, dictionary operations, \dots
                \end{itemize}   
                \end{block}     
\pause
\begin{block}{Definition (predefined)}
\begin{lstlisting}
inductive Option (α : Type) where
  | none : Option α
  | some : α → Option α

instance : Monad Option where
  pure x := Option.some x
  bind x f :=
    match x with
    | Option.none => Option.none
    | Option.some v => f v
\end{lstlisting}  
\end{block}
\end{frame}             

\begin{frame}[fragile]
  \frametitle{The List monad}
  \begin{block}{Useful for}
                \begin{itemize}         
                        \item Handling multiple results
                        \item Backtracking
                \end{itemize}   
                \end{block}     
\pause
\begin{block}{Definition (predefined)}
\begin{lstlisting}
def listBind : List α → (α → List β) → List β
  | List.nil, _ => List.nil
  | List.cons v vs, f => listAppend (f v) (listBind vs f)

instance : Monad List where
  pure x := List.cons x List.nil
  bind x f := listBind x f
\end{lstlisting}  
\end{block}
\end{frame}             

\begin{frame}[fragile]
  \frametitle{The IO Monad}
  \begin{block}{Required for}
  \begin{itemize}
        \item any kind of I/O
        \item side effecting operation
        \item implementation is machine dependent
  \end{itemize}
  \end{block}
\end{frame}

\begin{frame}
  \frametitle{Under the hood}
  Lean's \lstinline{Monad} class has more interesting structure:
  \begin{itemize}
  \item Pure, Bind
  \item Functors, Sequencing
  \end{itemize}
\end{frame}

\begin{frame}[fragile]
  \frametitle{Mapping and sequencing}
\begin{lstlisting}
class Functor (f : Type → Type) where
  map : (α → β) → f α → f β
\end{lstlisting}
  \begin{itemize}
  \item A \lstinline{Functor} is a type operator \lstinline{f} that has a \lstinline{map}
    operator.
  \item For lists, this is the well-known \lstinline{map}.
  \item Notation: \verb|<$>|
  \end{itemize}
  Any instance \lstinline{f} of the \lstinline{Functor} class must respect the \emph{functor laws:}
  \begin{itemize}
  \item preserve identify: \lstinline{id_map} --- \lstinline|∀ (x : α), map id x = x|
  \item preserve composition: \lstinline{comp_map} ---
    \lstinline{∀ (g : α → β) (h : β → γ) (x : f α), map (h ∘ g) x = map h (map g x)}
  \item cf.\ class \lstinline{LawfulFunctor}
  \end{itemize}
\end{frame}
\begin{frame}[fragile]
  \frametitle{Sequencing}
\begin{lstlisting}
class Seq (f : Type → Type) where
  seq : f (α → β) → (Unit → f α) → f β
\end{lstlisting}
  \begin{itemize}
  \item Function application lifted over \lstinline{f}.
  \item Notation: \lstinline{g <*> x}
  \end{itemize}
\begin{lstlisting}  
class SeqLeft (f : Type → Type) where
  seqLeft : f α → (Unit → f β) → f α

class SeqRight (f : Type → Type) where
  seqRight : f α → (Unit → f β) → f β
\end{lstlisting}
  
\end{frame}
\begin{frame}[fragile]
  \frametitle{Simple computations: Applicative}
\begin{lstlisting}
class Applicative (f : Type u → Type v)
    extends Functor f, Pure f, Seq f, SeqLeft f, SeqRight f where
  map g x := seq (pure g) (fun _ => x)
  seqLeft x y := seq (map (fun a _ => a) x) y
  seqRight x y := seq (map (fun _ b => b) x) y
\end{lstlisting}
  \begin{itemize}
  \item provides default implementations for \lstinline{map},
    \lstinline{seqLeft}, and \lstinline{SeqRight}.
  \item we can define an instance of \lstinline{Applicative} by implementing
    \lstinline{pure} and \lstinline{seq} --- explicit instances of the
    superclasses are not needed
  \item we'll be back, later
  \end{itemize}
\end{frame}

\begin{frame}[fragile]
  \frametitle{Real definition of the monad class}
\begin{lstlisting}
class Monad (m : Type u → Type v)
    extends Applicative m, Bind m where
  map f x := bind x (fun a => pure (f a))
  seq f x := bind f (fun g => Functor.map g (x ()))
  seqLeft x y := bind x (fun a => bind (y ()) (fun _ => pure a))
  seqRight x y := bind x (fun _ => y ())
\end{lstlisting}
  \begin{itemize}
  \item default implementations for \lstinline{map}, \lstinline{seq},
    \lstinline{seqLeft}, and \lstinline{seqRight}.
  \item Sufficient to implement \lstinline{pure} and \lstinline{bind} to obtain
    an instance of \lstinline{Monad}
  \end{itemize}
\end{frame}


\begin{frame}
  \frametitle{Challenges}
  \begin{itemize}
  \item what if there are multiple effects?
  \item need to combine monads
  \item sequence matters (e.g., exception and state)
  \item some monads do not combine at all
  % \item BUT we can go for something weaker
  \end{itemize}
  
\end{frame}


\end{document}


